\documentclass[11pt,a4paper]{article}
\usepackage[margin=2.3cm]{geometry}
\usepackage{fontspec}
\setmainfont{C059}[Path=/usr/share/fonts/urw-fonts/, Extension=.otf, UprightFont=*-Roman, BoldFont=*-Bold, ItalicFont=*-Italic, BoldItalicFont=*-BdIta]
\usepackage[english]{babel}
\usepackage{csquotes}
\usepackage{enumitem}
\usepackage{booktabs}
\usepackage{tabularx}
\newcolumntype{Y}{>{\raggedright\arraybackslash}X}
\usepackage{titlesec}
\usepackage{xurl}
\usepackage[hidelinks]{hyperref}
\setmonofont{NimbusMonoPS}[Path=/usr/share/fonts/urw-fonts/, Extension=.otf, UprightFont=*-Regular, BoldFont=*-Bold, Scale=MatchLowercase]
\emergencystretch=3em

\setlist[itemize]{leftmargin=1.4em, itemsep=3pt, topsep=4pt}
\titleformat{\section}{\normalfont\Large\bfseries}{\thesection.}{0.6em}{}
\titleformat{\subsection}{\normalfont\large\bfseries}{\thesubsection.}{0.6em}{}
\newcommand{\cit}[1]{{\small(#1)}}
\newcommand{\budget}[1]{\par\noindent\textit{Word budget for the essay: #1 words.}\par\smallskip}
\newcommand{\Braudel}{Braudel, \enquote{History and the Social Sciences}}
\newcommand{\Habib}{Habib, \enquote{Potentialities of Capitalistic Development}}

\begin{document}

\noindent{\LARGE\bfseries Seasonal Herding in the Himachal Hills, 1881--1961}\par
\smallskip
\noindent{\large Bullet points for the essay, by section}\par
\smallskip
\noindent 3 November 2026
\bigskip

\noindent Each bullet carries its source and PDF page. A
citation such as \cit{Brahmaur survey, PDF p.~55} names a 1961 village survey
monograph. The other sources are the 1881 census, Rose's \emph{Glossary},
the 1931 census, Newell's study of the Gaddis, the 1961 Punjab Census Atlas,
the Mahasu district handbook and the Weaving survey. Course readings are cited
by author, title and PDF page. Numbers are kept in the tables of
Appendix~\ref{app:data}. Figures read from page images are marked
\enquote{image-read}, and no person has checked them yet. The word budgets
add up to 3,800 words.

\tableofcontents
\bigskip

% ===================================================================
\section{Introduction}
\budget{250}

\subsection{Transhumance as the Phenomenon}
\begin{itemize}
\item Transhumance is the seasonal movement of herds between high summer
  pastures and low winter grazing by families who keep a fixed home and
  fields. Define it in these words at the first use.
\item Braudel names transhumance as his own example of a structure that holds
  people for centuries: \enquote{Man is a prisoner for long centuries of
  climates, of vegetations, of animal populations \dots{} Take the role of
  transhumance in mountain life.} \cit{\Braudel, PDF p.~10}
\item The Gaddis of the Ravi valley and the Kinnauras of the upper Sutlej move
  sheep and goats, and the Gujjars move buffaloes \cit{Brahmaur survey, PDF
  p.~55; Kanum survey, PDF p.~16; Maingal survey, PDF p.~52}.
\item An opening image is available in two descriptions of the same
  migration. In 1881 the Gaddis' wealth \enquote{consists of flocks of sheep
  and goats, which they feed half the year (the winter months) in the valleys
  of Kangra, and for the other half drive across the range into the
  territories of Chamba} \cit{1881 census, PDF p.~275}.
\item In 1961: \enquote{in the beginning of the winter (October and November),
  the flocks are driven to the low hills of Kangra, Nurpur, Pathankot and
  Jammu} \cit{Brahmaur survey, PDF p.~55}.
\end{itemize}

\subsection{Region}
\begin{itemize}
\item The study villages lie in four river systems. Brahmaur, Chitrari,
  Devi-Kothi and Maingal are in Chamba on the Ravi. Kanum, Kothi and Nachar are
  in Kinnaur on the Sutlej, and Shakrori lies lower on the Sutlej in Mahasu.
  Kupha (Pangi) and Tandi (Lahaul) are on the Chenab, and Bir is in Mandi on
  the Beas \cit{Appendix Table~\ref{tab:alt}}.
\item The survey villages are marked on the map \enquote{Himachal Pradesh,
  Village Surveys} \cit{Devi-Kothi survey, PDF p.~2} and on the Punjab map of
  villages selected for survey \cit{Tandi survey, PDF p.~5}.
\item The hills were divided among rulers. Chamba and Bushahr were princely
  states, Kangra was British territory, and Lahaul and Spiti lay in Punjab in
  1961 \cit{1881 census, PDF p.~275; Kanum survey, PDF p.~18; Tandi survey,
  PDF p.~5}.
\end{itemize}

\subsection{Period}
\begin{itemize}
\item The period runs from the 1881 census to the 1961 village surveys, with
  Rose's \emph{Glossary} (1911) and the 1931 census between them.
\item Within it the hill states became Himachal Pradesh, on 15 April 1948
  \cit{Kothi survey, PDF p.~16}, and the Forest Department took over the
  grazing tax \cit{Newell, PDF p.~19}.
\end{itemize}

\subsection{Research Question}
\begin{itemize}
\item How did the seasonal movement of flocks between high summer pastures and
  low winter grazing shape life in the Himachal hills between 1881 and 1961,
  and who controlled and profited from it?
\item As rule over the hills, control of the pastures, roads and trade
  changed, did the herding cycle persist as an adaptation to mountain and
  season, or did the ownership of flocks, pastures and shepherds' labour
  reshape it?
\item The introduction can close on the observation that both halves of the
  question turn out to be true at different speeds, without yet naming
  Braudel or Marx as the reason.
\end{itemize}

% ===================================================================
\section{Sources and Methodology}
\budget{500 (documents 100, computational methods 250, historical methods 150)}

\subsection{Documents}
\begin{itemize}
\item The 1961 Census village survey monographs of Himachal Pradesh, Vol.~XX,
  Part~VI\@. Tandi is in the Punjab series, Vol.~XIII, because Lahaul and Spiti
  lay in Punjab in 1961.
\item Four other 1961 volumes: Newell's study of the Gaddis in the
  \emph{Report on Scheduled Castes and Scheduled Tribes of Himachal Pradesh},
  the Weaving craft survey, the Mahasu district handbook, and the Punjab
  Census Atlas, Vol.~XIII, Part~IX.
\item The 1881 census of the Punjab, Results volume, which quotes the
  settlement officers Barnes and Lyall on the hill people \cit{1881 census,
  PDF p.~275}.
\item H.\,A. Rose, \emph{A Glossary of the Tribes and Castes of the Punjab and
  North-West Frontier Province} (1911 to 1919).
\item The 1931 census of the Punjab, Part~I, Report.
\item Page numbers in the essay are PDF page numbers, which are the page
  counter of a PDF viewer. State this once.
\end{itemize}

\subsection{Computational Methods}
\begin{itemize}
\item Every page of every volume was scored by the frequency of herding words,
  and the highest-scoring pages were selected \cit{Appendix
  Table~\ref{tab:pipe}}.
\item A language model read each selected page and proposed claims. A claim
  is one statement from one page, with a subject, a relation, an object, any
  quantity, months and place, and a verbatim quotation.
\item A script kept a claim only if its quotation is found on the cited page.
  A second model pass judged each claim against its page, and claims judged
  not supported were left out.
\item The livestock tables and Newell's fee table were read from page images,
  because their OCR scrambles the columns. Say that no person has checked these
  figures, unless the group checks them before submission.
\item The claims form an evidence graph, which is a network whose edges are
  sourced claims. Every table, chart and map of the essay is a query on it.
\item Each claim carries a tag for its conjuncture: the cycle, two sovereigns,
  the raja's pastures, the Forest Department, roads and wages, trade and credit,
  or the standing social order. The tag makes it possible to count the
  evidence for each conjuncture by the year of its source \cit{Appendix
  Table~\ref{tab:graph}}.
\item Payments in the graph are turned into flows from the party who gives to
  the party who receives, so that the receivers can be compared across
  conjunctures \cit{Appendix Table~\ref{tab:flows}}.
\item Keyword rates per 10,000 words were measured for each source layer and
  each survey, as a distant reading of what each source talks about
  \cit{Appendix Tables~\ref{tab:trade} and~\ref{tab:layers}}.
\item Altitudes and households were taken from each survey's own statement and
  checked against SRTM satellite elevation. Most surveys agree with it within
  300 m, and Nachar, Chaupal and Rajana do not \cit{project analysis A18}.
\item A text-reuse search compared every page with every page of the earlier
  sources. Barnes on the Gujars in 1881 reappears in Rose, so the two pages are
  one witness \cit{1881 census, PDF p.~267; Rose, vol.~2, PDF p.~346}.
\item Newell's report repeats Rose on its later pages. Newell's pages on
  grazing, fees, permits and the migration hold no copied passage
  \cit{Appendix Table~\ref{tab:pipe}}.
\item Bloch supplies the rule behind the reuse step: two witnesses with the
  same details mean \enquote{that one of them had copied the other or that both
  had copied a common model} \cit{Bloch, \emph{The Historian's Craft}, PDF
  p.~36}.
\item Bloch also supplies the caution for the model-proposed claims:
  \enquote{There is no reliable witness in the absolute sense. There is only
  more or less reliable testimony.} \cit{Bloch, \emph{The Historian's Craft},
  PDF p.~23}
\end{itemize}

\subsection{Historical Methods}
\begin{itemize}
\item Braudel divides historical time into the longue durée of geography and
  structures, the conjuncture of medium-term cycles, and the event \cit{\Braudel,
  PDF pp.~5--10}.
\item The three levels run at the same time, at different speeds. Braudel:
  \enquote{The longue durée, cyclical phases (conjoncture), and events fit
  together easily, for they all are measurements on the same scale.}
  \cit{\Braudel, PDF p.~29}. Call this temporal parallelism and define it as
  the simultaneous running of the three levels.
\item A conjuncture is measured by series, such as prices, wages and interest
  rates \cit{\Braudel, PDF p.~7}. The fee series in Newell is such a series.
\item Bloch's comparative method sets side by side societies that share
  conditions and differ in one respect: \enquote{At the bottom of nearly all
  criticism there is a problem of comparison} \cit{Bloch, \emph{The Historian's
  Craft}, PDF p.~32}.
\item The Marxist method identifies the material basis of a social formation,
  meaning its means and relations of production, and its dialectical
  contradictions \cit{Alam, \enquote{Marxist History} slides, PDF pp.~10, 15}.
\item State the plan of the essay in one sentence: Braudel's levels order the
  evidence in time, and Marx's categories explain the relations inside each
  level. Leave the interpretation for Section~\ref{sec:interp}.
\end{itemize}

% ===================================================================
\section{Geography}
\budget{500}

\subsection{Time}
\begin{itemize}
\item Braudel's geographical time is the slowest level, the time of
  \enquote{man in his relationship to the environment}. Its structures last
  for centuries: \enquote{a reality that time can only slowly erode, one that
  goes on for a long time} \cit{\Braudel, PDF p.~9}.
\item Structures both support and limit: \enquote{all structures are
  simultaneously pillars and obstacles \dots{} Think of how difficult it is to
  break through certain geographical frameworks} \cit{\Braudel, PDF p.~9}.
\item Routes belong to the same level: \enquote{Look at the endurance of roads
  and trade routes} \cit{\Braudel, PDF p.~10}.
\item Le Roy Ladurie sets the limit of this level: \enquote{The aim of
  climatic history is not to explain human history} \cit{Le Roy Ladurie,
  \enquote{The History of Rain and Fine Weather}, PDF p.~3}. Use this sentence
  to prepare the reader for what geography cannot explain.
\item The calendar evidence belongs to Section~\ref{sec:movement}. This
  subsection introduces the concept only.
\end{itemize}

\subsection{Altitude and River Valleys}
\begin{itemize}
\item The Atlas divides the region by height. North of the low hills lies a
  zone of high hills in Kangra, and north of that the mountainous areas.
  Lahaul and Spiti, separated from Kangra by high ranges, is very high
  \cit{Punjab Census Atlas, PDF p.~317, notes to Map~4; map on PDF p.~24}.
\item Barnes placed the Gaddis in a band of cultivation, with \enquote{little or
  no cultivation} above it \cit{1881 census, PDF p.~275}.
\item The village surveys state their own heights, from Shakrori \enquote{at a
  height of 1,700 ft.} to Tandi \enquote{on a ten thousand feet high flat
  ground} \cit{Shakrori survey, PDF p.~11; Tandi survey, PDF p.~18}.
\item In the Ravi, Sutlej and Chenab valleys, households in every village
  above a threshold of altitude keep large flocks, and households in every
  village below it keep few \cit{Appendix Table~\ref{tab:alt}, image-read}.
\item The rank correlation between altitude and animals per household across
  the survey villages is positive, and the correlation between altitude and the
  landholding castes' share is close to zero \cit{Appendix
  Table~\ref{tab:alt}}. Height set the size of flocks, not who owned them.
\item Chaupal stands as high as the herding villages, at \enquote{7,685 feet},
  but on the Yamuna side, and keeps few animals \cit{Chaupal survey, PDF
  p.~12; Appendix Table~\ref{tab:alt}}.
\item The Atlas names four perennial rivers, the Jumna, the Sutlej, the Beas
  and the Ravi \cit{Punjab Census Atlas, PDF p.~317, notes to Map~5; map on
  PDF p.~25}.
\item Brahmaur lies on \enquote{the Budhal stream which joins Ravi}, and Kanum
  lies above \enquote{the River Sutlej} \cit{Brahmaur survey, PDF p.~11; Kanum
  survey, PDF p.~15}.
\item Newell describes the Gaddi country as a chain of river units, each
  \enquote{a 40 mile length of river between high mountains} \cit{Newell,
  PDF p.~10}.
\item The village notional maps show each village in its valley
  \cit{Brahmaur survey, PDF p.~2; Kanum survey, PDF p.~2; Tandi survey,
  PDF p.~17}.
\end{itemize}

\subsection{Vegetation and Rainfall}
\begin{itemize}
\item Rainfall in Punjab rises from the south-west towards the north-east and
  the hills \cit{Punjab Census Atlas, PDF p.~318, notes to Map~8; map on PDF
  p.~28}.
\item In the mountains, \enquote{temperature, altitude, aspect and slope} add
  to rainfall in setting the vegetation. In Lahaul and Spiti the rainfall is
  low, vegetation is confined to the river valleys, and elsewhere there is bare
  rock \cit{Punjab Census Atlas, PDF p.~320, notes to Map~13; map on PDF
  p.~44}.
\item Forests occupy the largest area of Chitrari, and the villagers take
  firewood, timber and grass from them \cit{Chitrari survey, PDF p.~38}.
\item The people of Brahmaur take timber from the Brahmni reserved forest
  through the Forest Department \cit{Brahmaur survey, PDF p.~24}.
\item On the route to the plains, the Gaddis use \enquote{pasturage amid the
  forests lower down and pasturage on the plains} \cit{Newell, PDF p.~18}.
\item Sheep \enquote{live mainly on grazing}, so they are numerous where there
  is fallow or waste land \cit{Punjab Census Atlas, PDF p.~367}. The Atlas
  sheep and goat maps place them in Kangra, Simla, and Lahaul and Spiti
  \cit{Punjab Census Atlas, PDF pp.~290, 291}.
\end{itemize}

\subsection{Climate and Seasons}
\begin{itemize}
\item Snow \enquote{usually begins to fall on the passes in October and closes
  them by the middle of November}, and \enquote{the passes re-open in May}
  \cit{Kupha survey, PDF p.~111}.
\item The pass from Kugti to Lahaul is \enquote{a bare mountain top at a height
  of over 16,000} feet, and it is rare for anyone with animals \enquote{to take
  less than three days actually on the tops of the mountain} \cit{Newell, PDF
  p.~18, image-read}.
\item The grazing near the village is not enough: \enquote{these tarkals are
  not sufficient for the upkeep of the entire sheep and goat wealth of the
  village and people have therefore to migrate to other pastures}
  \cit{Brahmaur survey, PDF p.~55}.
\item In Tandi, spinning is done indoors in winter and weaving in summer, so
  the seasons order household work as well as herding \cit{Tandi survey, PDF
  p.~66}.
\item The section can end on the point that the seasons fixed when the flocks
  could move, but not who owned them.
\end{itemize}

% ===================================================================
\section{Socio-economy}
\budget{750 (movement 200, flocks 250, trade and commerce 150, social structure 150)}

\subsection{Movement: Geography and Transport}\label{sec:movement}
\begin{itemize}
\item In 1881 the Gaddis spend \enquote{half the year (the winter months) in the
  valleys of Kangra} \cit{1881 census, PDF p.~275}.
\item In 1911, \enquote{the yearly exodus to Kangra takes place in October and
  November, and the return journey in April and May}, with summers \enquote{in
  the higher mountains of Pangi and Lahul} \cit{Rose, vol.~2, PDF p.~297}.
\item In 1931, \enquote{the hillmen of the lofty ranges leave their homes with
  their sheep and cattle to spend the winter in the valleys and plains}
  \cit{1931 census, PDF p.~142}.
\item In 1961 the Gaddis \enquote{visit Pangi with their flocks in July and
  remain in the valley till September} \cit{Kupha survey, PDF p.~113}, and of
  Brahmaur's households \enquote{as many as 61 move out in winter months}
  \cit{Brahmaur survey, PDF p.~63}.
\item The 1881 and 1931 censuses record a winter descent without naming months.
  Only Rose and the 1961 surveys name the months, and they agree
  \cit{Appendix Table~\ref{tab:cal}}.
\item The Gujjars of Maingal take their buffaloes up in May and return in
  September \cit{Maingal survey, PDF p.~52}. The Kinnaur flocks spend
  \enquote{the chilly months} in the lower hills \enquote{till the end of April}
  \cit{Kanum survey, PDF p.~85}.
\item The flocks cross from one river system to another. Brahmaur's flocks go
  south over the Dhauladhar to Kangra in winter and north to Pangi and Lahaul
  in summer \cit{Brahmaur survey, PDF p.~55}.
\item On the route, families camp at known rock shelters, sometimes many
  families under one rock \cit{Newell, PDF p.~36}.
\item Transport changed by 1961. The route to western Tibet over the Rohtang
  Pass, \enquote{previously only a bridle path}, was made jeepable from Manali
  to Keylong \cit{Tandi survey, PDF p.~57}.
\item A Gaddi shepherd in Chitrari says that on \enquote{our beaten track for
  centuries a jeep has started plying} \cit{Chitrari survey, PDF p.~24,
  image-read}. A jeepable road reached Kanum \cit{Kanum survey, PDF p.~133}.
\item The Atlas records that Kangra and Lahaul and Spiti \enquote{were generally
  without metalled roads} in 1947 and had better road connections by 1961
  \cit{Punjab Census Atlas, PDF p.~372; maps on PDF pp.~308, 309}.
\end{itemize}

\subsection{Flocks: Ownership, Labour and Grazing Rights}
\begin{itemize}
\item Landholding castes owned most of the sheep and goats in every caste table
  \cit{Appendix Tables~\ref{tab:class} and~\ref{tab:own}, image-read}.
\item In Brahmaur, Brahmans and Rajputs are most of the households, and a part
  of them own almost all the animals. The Sipis own few, and the one Rehara
  household owns none \cit{Brahmaur survey, PDF p.~54; Appendix
  Table~\ref{tab:brahmaur}}.
\item In Nachar a family keeps \enquote{100 to 200 sheep and goats}, while
  \enquote{the Lambardar of the village has about 1,100 strong flock}. The
  headman was also the largest owner \cit{Nachar survey, PDF p.~33}.
\item In Rajana, Pangna and Chaunri, where households kept a few goats, Koli
  households held most of the animals \cit{Appendix Table~\ref{tab:own}}.
\item In Brahmaur most families herd their own flocks. A malundi is \enquote{a
  shepherd who tends his own flock}, and \enquote{pohals are hired in summer
  season by a very few people in rare cases} \cit{Brahmaur survey, PDF
  p.~55}.
\item A pohal is paid in kind and a little cash: \enquote{free meals},
  \enquote{free clothing and wool for the family}, \enquote{two or four sheep
  per hundred sheep and goats} and \enquote{Rs. 2/- in cash per month}
  \cit{Brahmaur survey, PDF p.~55}.
\item A Kanum trader \enquote{had four fawals-shepherds who were paid Rs. 60
  each (annually) free meals, 8 kilogram of wool, shoes and clothes}
  \cit{Kanum survey, PDF p.~132}.
\item In Tandi every household pays the village shepherd \enquote{one patha
  \dots{} of buckwheat per sheep or goat for the season} \cit{Tandi survey,
  PDF p.~65}.
\item Devi-Kothi farmers \enquote{may also send their sheep and goats with
  professional graziers to high hills} \cit{Devi-Kothi survey, PDF p.~48}.
\item Newell: \enquote{the majority of Gaddis are not shepherds} \cit{Newell,
  PDF p.~37}. In Chitrari \enquote{the persons who take to this seasonal
  migration are the ones who have a sizeable flock to tend} \cit{Chitrari
  survey, PDF p.~23}.
\item Under princely rule, \enquote{In Chamba all grazing rights belonged to the
  rajah of Chamba}, and the right to collect fees \enquote{was auctioned by the
  rajah of Chamba to the highest bidder, who was usually a Gaddi} \cit{Newell,
  PDF p.~18}.
\item Near Brahmaur the ruler leased the pastures to a contractor, who allotted
  them to shepherds on condition that the flocks manure the fallow fields at
  night \cit{Brahmaur survey, PDF p.~55}.
\item In Kanum, by contrast, \enquote{every villager has a right of utilizing the
  pastures} under the Bushahr ruler as Primary Owner \cit{Kanum survey, PDF
  pp.~18, 133}.
\item After 1948 \enquote{the pastures are allotted direct to the shepherds by the
  Forest Department}, and each shepherd has a fixed dhar for \enquote{a fixed
  number of sheep and goats} \cit{Brahmaur survey, PDF p.~55}.
\item The tax receipt \enquote{is however not issued unless the shepherd has a
  permit issued by the Forest Department} \cit{Newell, PDF p.~19}, and
  \enquote{the Forest Department do not now grant any additional permit}
  \cit{Brahmaur survey, PDF p.~46}.
\item In Devi-Kothi \enquote{A distinction is made between right holder graziers
  and others in payment of fee and penalty} \cit{Devi-Kothi survey, PDF
  p.~47}. The rights were \enquote{given to them by the early rulers of Chamba}
  \cit{Devi-Kothi survey, PDF p.~49}.
\item In the survey's Table XIV the right holders pay less only for buffaloes,
  in both fee and penalty. Sheep, goats and cattle pay the same whether or not
  the grazier holds rights \cit{Devi-Kothi survey, PDF p.~47, image-read;
  Appendix Table~\ref{tab:fees}}.
\item Do not use the Devi-Kothi sentence \enquote{the fee is Rs. 5/- from
  non-right holders and Re. 1/- from a right holder} as a grazing fee. It is
  the fee for a permit to extract dhup, a herb, from the hill tops
  \cit{Devi-Kothi survey, PDF p.~48}. The old draft essay used it as a grazing
  fee.
\item The raja's fee and the Department's fee are in different units. The raja's
  contractor charged per head for each grazing district, and the districts
  \enquote{were so arranged that shepherds might wander from one to the other}
  \cit{Newell, PDF p.~18}. The Department in Brahmaur charged per hundred heads
  for the summer season, and \enquote{A shepherd has to pay the grazing tax twice
  a year} \cit{Brahmaur survey, PDF p.~55}. A flock's total payment under the
  raja depended on how many districts it grazed, which no source records, so
  the two fees cannot be compared per head per year \cit{Appendix
  Table~\ref{tab:fees}}.
\item Newell's 1962 rates give one comparison in the same units, across three
  areas under the Department. In Lahul and Kangra the Department charged more
  for a goat than for a sheep, and in Brahmaur it charged both the same \cit{Newell, PDF p.~19,
  image-read; Appendix Table~\ref{tab:fees}}.
\item The number of fee-paying sheep and goats fell in one step and did not
  recover \cit{Newell, PDF p.~38, image-read; Appendix
  Table~\ref{tab:series}}.
\item Permits became property: \enquote{these permits are freely bought and sold
  amongst shepherds so that shepherds have now become closed corporations}
  \cit{Newell, PDF p.~37}.
\end{itemize}

\subsection{Trade and Commerce}
\begin{itemize}
\item The Gaddis' woollens \enquote{command a ready market in the low hills where
  people sell them during their stay there in winter months} \cit{Brahmaur
  survey, PDF p.~59}.
\item Kanum's people went down with \enquote{their flocks laden with chilgozas,
  wool, woollen patties, nuts, pasham, jira and other local herbs} and sold
  them \enquote{either in Lavi fair at Rampur} or in the lower villages
  \cit{Kanum survey, PDF p.~16}.
\item From Tibet \enquote{they brought wool, pasham, butter, ghee, goat, sheep
  and salt} \cit{Kanum survey, PDF p.~131}, and a trader's journey
  \enquote{began in July and he was back in Kanum usually by November}
  \cit{Kanum survey, PDF p.~132}.
\item The Lavi fair at Rampur, in the second week of November, is \enquote{the
  biggest trade fair in Himachal} \cit{Mahasu handbook, PDF p.~56; fairs map on
  PDF p.~4}.
\item Rose records Tibetan \enquote{trade sheep} crossing the Baralacha pass
  \cit{Rose, vol.~1, PDF p.~523}.
\item Trade words are far more frequent in the Sutlej surveys than in the Ravi
  surveys \cit{Appendix Table~\ref{tab:trade}}.
\item The low villages bought from passing herders. Bir's animals were
  \enquote{purchased from Gaddis who wander near about in winter, for 10 to 12
  rupees}. Batal bought wool \enquote{from the Kinnauras who stay near about the
  village} \cit{Bir survey, PDF p.~43; Batal survey, PDF p.~46}.
\item After the trade with Tibet stopped, \enquote{an acute shortage of wool and
  Pasham is being experienced} \cit{Weaving survey, PDF p.~16}, and Chergaon's
  buyers turned from the Lavi fair to Dodra Kawar \cit{Chergaon survey, PDF
  p.~61}.
\item In the graph, interest paid by households to moneylenders is the most
  frequent relation of payment in the whole corpus \cit{Appendix
  Table~\ref{tab:flows}}.
\item Co-operative societies appear as lenders in the 1961 surveys of Kalath,
  Pangna, Kot, Brahmaur and Rajana. They are a new creditor beside the
  moneylender and the shopkeeper \cit{Appendix Table~\ref{tab:flows}}.
\item Credit ran alongside trade. In Maingal the Gujjars \enquote{borrow money at
  30\% annual interest} from lenders who also buy their milk products
  \cit{Maingal survey, PDF p.~42}.
\item The Gujjars sell their ghee \enquote{to a particular shopkeeper at Chamba
  who pays them much less than the market price} \cit{Maingal survey, PDF
  p.~60}, and those in debt \enquote{are obliged to sell their produce to them}
  \cit{Maingal survey, PDF p.~61}.
\item In Devi-Kothi the shopkeeper from Kangra charges high rates, and Bhot
  traders from Pangi sell goods for grain \cit{Devi-Kothi survey, PDF p.~53}.
\end{itemize}

\subsection{Social Structure: Family Units, Caste, Customs and Traditions}
\begin{itemize}
\item The household is the unit of migration. On the winter route \enquote{the
  wife usually accompanies her husband}, cooks and looks after the children,
  while the husband looks after the sheep \cit{Newell, PDF p.~36}.
\item \enquote{Even during migration, she carries an equal load on her back and
  contributes equally to the family earnings} \cit{Brahmaur survey, PDF
  p.~79}.
\item In Nachar, households with large flocks spare one or two members as
  shepherds \cit{Nachar survey, PDF p.~33}.
\item Occupations follow caste. In Maingal, \enquote{Occupations are classified
  better on caste basis}: Gujjars rear buffaloes, and Gaddi Brahmins and
  Rajputs farm and keep sheep and goats \cit{Maingal survey, PDF p.~38}.
\item Debt runs across caste lines: \enquote{Gujjars had to pay food grain debts
  to Gaddi brahmins} \cit{Maingal survey, PDF p.~38}.
\item In Kanum the landlord who is also the moneylender has lent to
  \enquote{the scheduled castes and others} \cit{Kanum survey, PDF p.~53}.
\item Custom shapes the economy of the village. No poultry is kept in Brahmaur
  proper, because the birds would \enquote{defile the sacred Chaurasi area}
  \cit{Brahmaur survey, PDF p.~54}.
\item In Kanum the temple lends to every household and collects interest at the
  fair of Kangyur \cit{Kanum survey, PDF p.~82}.
\item \enquote{At the top of every pass is a shrine at which the herdsmen
  sacrifice a beast or the traveller presents a small coin} \cit{Newell, PDF
  p.~19}. The route of the flocks is also a sacred geography.
\item In Tandi one household keeps the village yak in winter, and the whole
  village shares the cost \cit{Tandi survey, PDF p.~65}.
\item Bushahr's subjects owed the ruler labour service, \enquote{including
  `Athada Begar'} \cit{Kanum survey, PDF p.~18}.
\item The 1931 Punjab Social Map shows Chamba and Bashahr as separate states
  and draws the Depressed Classes as a category \cit{1931 census, PDF p.~24}.
\end{itemize}

% ===================================================================
\section{Political Events}
\budget{300}

\begin{itemize}
\item In 1874 \enquote{an attempt was made to lease the dhars or bigger grazing
  grounds} of Pangi, and \enquote{it provoked much opposition and was finally
  abandoned} \cit{Kupha survey, PDF p.~113}.
\item In 1881 the Gaddis \enquote{hold lands on this side and also in Chamba}
  and, \enquote{whenever fined by the Kangra authorities}, paid \enquote{a
  similar penalty into the Chamba treasury}. Barnes adds that \enquote{our
  institutions have taught them greater independence} \cit{1881 census, PDF
  p.~275}.
\item Gaddis crossing into Lahul \enquote{would need to pay a grazing tax to
  both the Chamba and Lahuli ruler} \cit{Newell, PDF p.~18}.
\item On 15 April 1948 the hill states became Himachal Pradesh, and Kinnaur
  left the state of Rampur Bushahr \cit{Kothi survey, PDF p.~16}.
\item After the merger, \enquote{these rights to grazing tax have now been taken
  over by the Government of India State Forest Department} \cit{Newell, PDF
  p.~19}, and \enquote{today the power of the forest department is paramount}
  \cit{Newell, PDF p.~15}.
\item When the Indian Union incorporated Chamba, it introduced \enquote{the idea
  of private property in land}, and \enquote{in 1953 a Settlement was instituted
  in Brahmaur tahsil} \cit{Newell, PDF p.~19}.
\item Shipki La \enquote{was open to all traders without passport} before
  \enquote{the India-China dispute} \cit{Kanum survey, PDF p.~51}, and the Gartok
  trade \enquote{has now practically come to an end} \cit{Kanum survey, PDF
  p.~16}. No source in the corpus dates the closing.
\item Newell marks one year of the fee series as \enquote{the year after the
  sheep were struck by disease} \cit{Newell, PDF p.~38}.
\item Keep this section to a list of events in order. Their meaning belongs to
  Section~\ref{sec:interp}.
\end{itemize}

% ===================================================================
\section{Interpretation}\label{sec:interp}
\budget{1,300}

\noindent The hypothesis is that Braudel's levels give the frame of the
history, and Marx's categories explain the relations and contradictions inside
that frame. The bullets below let the reader reach this through the cases. The
essay can show it in each subsection and name it once, near the end.

\subsection{Three Clocks over One Valley}
\begin{itemize}
\item Read the evidence on Braudel's three levels at once. The calendar and the
  altitude threshold are the longue durée. The raja's fee contract, the Forest
  Department's permits, the Tibet trade and road building are conjunctures. The
  lease attempt of 1874, the merger of 1948, the India-China dispute and the
  sheep disease are events \cit{Sections~3 to~5 of the essay}.
\item The three levels run together: \enquote{The longue durée, cyclical phases
  (conjoncture), and events fit together easily} \cit{\Braudel, PDF p.~29}.
\item The Brahmaur surveyor wrote both clocks into one sentence: roads offer
  \enquote{new avenues of earning livelihood}, which \enquote{is likely to change
  the cycle of their life followed since ages} \cit{Brahmaur survey, PDF
  p.~87}.
\item The frame answers the first half of the research question. The months,
  the routes and the threshold of altitude did not change between 1881 and
  1961 \cit{Appendix Tables~\ref{tab:alt} and~\ref{tab:cal}}.
\item The frame leaves questions open, and these questions lead to Marx. Why
  did Chaupal, as high as Brahmaur, keep few animals? Why did the Gaddis of
  Maingal settle? Why did fee-paying flocks fall and stay down?
\end{itemize}

\subsection{Limits of the Geographical Frame}
\begin{itemize}
\item Altitude made large flocks possible, but whether a family migrated
  depended on whether it owned a flock large enough to make the journey worth
  it \cit{Chitrari survey, PDF p.~23}.
\item Within the same valley and under the same ruler, Brahmaur families went
  with their own flocks, Devi-Kothi farmers paid graziers, and Maingal's Gaddis
  farmed \cit{Brahmaur survey, PDF p.~55; Devi-Kothi survey, PDF p.~48; Maingal
  survey, PDF p.~52}. Bloch's comparison rules out geography as the cause of
  this difference.
\item The same geography carried different relations of production. The owner
  who worked his own flock, the hired pohal paid in kind, the grazier paid per
  head and the village shepherd paid in grain are four relations on one
  mountain \cit{Brahmaur survey, PDF p.~55; Devi-Kothi survey, PDF p.~48; Tandi
  survey, PDF p.~65}.
\item Marx's first step applies here: \enquote{Identify the `Material Basis' of
  social formation} \cit{Alam, \enquote{Marxist History} slides, PDF p.~10}. The
  material basis is the flocks and the pastures, and who held them.
\item Braudel himself credits Marx with \enquote{real social models, based on the
  historical longue durée} \cit{\Braudel, PDF p.~32}. Quote this to show that
  the two lenses were not opposed even for Braudel.
\end{itemize}

\subsection{Each Conjuncture Has Its Material Basis and Its Contradiction}
\begin{itemize}
\item Two sovereigns, 1881. One people stood under two claims, British Kangra
  and Chamba. Barnes saw colonial rule weakening the raja's side of the claim
  \cit{1881 census, PDF p.~275; Newell, PDF p.~18}.
\item The raja's pastures, to about 1950. The raja owned the pastures and the
  herders used them. A share of the product was taken as a cash fee, as labour
  under the Athada Begar, as wool in place of revenue, and as manure under the
  contract \cit{Newell, PDF pp.~14, 18; Kanum survey, PDF p.~18; Brahmaur
  survey, PDF p.~55}.
\item Marx counts the state's claim on the product as its department of
  \enquote{Finance, or the plunder of the interior} \cit{Marx, \enquote{The
  British Rule in India}, PDF p.~2}. Habib separates such revenue from rent:
  \enquote{a tax on the crop} \cit{\Habib, PDF p.~8}. The raja's grazing fee
  was both, because he was ruler and owner at once.
\item Newell describes the princely order in terms close to Marx's: \enquote{the
  whole of Chamba State was nationalised, that is belonged to the nation in the
  person of the rajah}. The raja had to allocate land \enquote{on certain
  customary principles}, and land \enquote{could thus not be used as security for
  debt} \cit{Newell, PDF p.~19}.
\item After incorporation into Himachal Pradesh the state introduced
  \enquote{the idea of private property in land}, with a Settlement of Brahmaur
  tahsil in 1953 \cit{Newell, PDF p.~19}. Marx saw the zamindari and ryotwari
  settlements of British India as \enquote{two distinct forms of private property
  in land} \cit{Marx, \enquote{The Future Results of British Rule in India}, PDF
  p.~1}. The Settlement of 1953 is the same transformation, arriving in Chamba
  after independence.
\item Under the raja, land could not secure a debt. Once land became private
  property, it could. The Brahmaur and Devi-Kothi surveys each have a section on
  mortgages \cit{Brahmaur survey, PDF p.~3; Devi-Kothi survey, PDF p.~5, both
  contents pages}. Read those sections before the essay links the Settlement to
  mortgaged land.
\item The fee contractor was usually a Gaddi, so the raja's claim was collected
  through the herding community itself \cit{Newell, PDF p.~18}.
\item The contradiction was ownership against customary use. The 1874 lease was
  abandoned after opposition, so custom set a limit on the ruler's claim
  \cit{Kupha survey, PDF p.~113}. The Manifesto's image fits: old relations
  \enquote{became so many fetters} \cit{Marx and Engels, \emph{Manifesto},
  PDF p.~4}.
\item The Forest Department, from 1948. The state became the owner of the
  pastures and collected the fee through its own employee, without the Gaddi
  contractor \cit{Newell, PDF p.~19}. Aymard describes such \enquote{elimination
  of middlemen} in the collection of a tax \cit{Aymard, \enquote{The Cost of
  War}, PDF p.~3}.
\item The unit of the charge changed. The raja's contractor collected per head
  for each grazing district a flock entered, and the Department collected per
  head for a season, under a permit for a fixed number of animals on a fixed
  dhar \cit{Newell, PDF p.~18; Brahmaur survey, PDF pp.~46, 55; Appendix
  Table~\ref{tab:fees}}. The charge was tied to the district under the raja and
  to the permit holder under the Department.
\item The contradiction was conservation against subsistence. Erosion led to
  passes and closures, and closures led to resentment \cit{Chitrari survey,
  PDF p.~23; Kupha survey, PDF p.~136}. Kosambi: \enquote{Fusion and tolerance
  become impossible when \dots{} there is not enough of the surplus product to
  go around} \cit{Kosambi, PDF p.~21}.
\item The resolution was a new form of property. A capped permit became
  a commodity that shepherds bought and sold, and the holders became \enquote{closed
  corporations} \cit{Newell, PDF p.~37}. Habib describes the same passage for
  zamindari rights: \enquote{owing to monetization, the zamindari right became a
  fully salable commodity} \cit{\Habib, PDF p.~15}.
\item The slides' rule of quantity and quality applies: \enquote{A sufficiently
  large change in quantity changes the nature of the thing} \cit{Alam,
  \enquote{Marxist History} slides, PDF p.~8}. A limit on the number of permits
  turned a licence into property.
\item Old privilege survived inside the new regime. The Department kept the
  distinction between right holders, whose rights came from the rajas, and
  others, and charged right holders less for buffaloes \cit{Devi-Kothi survey,
  PDF pp.~47, 49, image-read}. The slides: \enquote{Colonial Relations Can Continue Even After
  Political Independence} \cit{Alam, \enquote{Marxist History} slides, PDF
  p.~17}.
\item Roads and wages, the 1950s. The pohal was paid in meals, clothing, wool
  and a few animals, and the road labourer in cash by the day \cit{Brahmaur
  survey, PDF pp.~45, 55; Appendix Table~\ref{tab:wages}}. In Kanum, labourers \enquote{prefer to work on the
  roads rather than work as `fawals'} because \enquote{they get better wages}
  \cit{Kanum survey, PDF p.~51}.
\item The contradiction was between flock owners who needed hired shepherds and
  shepherds who could now be paid higher wages on the roads. The Manifesto names the new tie
  as \enquote{callous `cash payment'} \cit{Marx and Engels, \emph{Manifesto}, PDF
  p.~2}.
\item Marx expected roads to break the isolation of Indian villages:
  \enquote{the absence of roads perpetuated the village isolation} \cit{Marx,
  \enquote{The Future Results of British Rule in India}, PDF p.~2}. The jeep on
  the shepherds' track is the local case \cit{Chitrari survey, PDF p.~24}.
\item Credit ran through every conjuncture. The Gujjars owned buffaloes and were
  all in debt, so their surplus was taken through interest and the price of
  ghee, without any claim on their animals \cit{Maingal survey, PDF pp.~42,
  48, 60}.
\item Habib names the beneficiaries of such debt as headmen, \enquote{professional
  usurers} and \enquote{merchants (who often gave advances to the peasants to
  establish claims on their produce)} \cit{\Habib, PDF p.~14}. The Manifesto
  lists \enquote{the landlord, the shopkeeper, the pawnbroker} \cit{Marx and
  Engels, \emph{Manifesto}, PDF p.~5}. Kanum's landlord-moneylender and the
  Chamba ghee buyer are both on that list \cit{Kanum survey, PDF p.~53;
  Maingal survey, PDF p.~60}.
\end{itemize}

\subsection{Evidence from the Knowledge Graph}
\begin{itemize}
\item The main receiver of payments changes with the conjuncture. In the
  raja's conjuncture it is the Raja of Chamba and the fee contractor, after
  1948 the Forest Department, and in trade and credit the moneylender and the
  shopkeeper \cit{Appendix Table~\ref{tab:flows}; flow chart in
  \path{analysis/surplus/flows.png}}. This is the Marxist question of who
  takes the surplus, answered by a count over the whole corpus.
\item The evidence for the raja's conjuncture comes mostly from Newell, writing
  in 1961 about the years before, and from Rose. The 1931 census has no claim
  on it \cit{Appendix Table~\ref{tab:graph}}. The princely order is therefore
  seen largely through a later witness, which Bloch's criticism asks the essay
  to say \cit{Bloch, \emph{The Historian's Craft}, PDF p.~23}.
\item The Forest Department's conjuncture is attested only in the 1961
  sources, and words of forest and state are absent from the 1881 and Rose
  layers \cit{Appendix Tables~\ref{tab:graph} and~\ref{tab:layers}}. The
  conjuncture appears in the vocabulary of the sources as well as in their
  content.
\item Words of wage work and debt are more frequent in the 1931 census than in
  the 1881 census \cit{Appendix Table~\ref{tab:layers}}. The colonial rates
  are over whole volumes, not only the hills, so they show a direction and not
  a measure for the hills.
\item The conjuncture of two sovereigns is attested by very few claims \cit{Appendix
  Table~\ref{tab:graph}}. The essay can present it as an observation of Barnes
  and Newell, not as a measured pattern.
\item The calendar holds in independent witnesses. Barnes's passage on the
  Gujars was copied into Rose, but Rose's account of the Gaddi year and the
  1961 accounts are not copies \cit{project analysis A15, A19}. Bloch:
  \enquote{agreement in the errors of the data \dots{} seems inconceivable to
  us} \cit{Bloch, \emph{The Historian's Craft}, PDF p.~42}.
\end{itemize}

\subsection{Events as the Points Where the Levels Meet}
\begin{itemize}
\item Braudel: \enquote{An event is an explosion, something that has `the sound of
  newness'} \cit{\Braudel, PDF p.~5}. Each event in the sources shows where a
  structure and a contradiction met.
\item The lease of 1874 is an event in which the contradiction between the
  ruler's ownership and the herders' use came into the open. Its failure kept
  customary grazing in place \cit{Kupha survey, PDF p.~113}.
\item The merger of 1948 changed the sovereign but not the migration. Marx,
  quoting a Commons report on Indian villages: \enquote{while the village remains
  entire, they care not to what power it is transferred} \cit{Marx, \enquote{The
  British Rule in India}, PDF p.~3}. What the merger did change was the form of
  the claim on the pastures \cit{Newell, PDF p.~19}.
\item The India-China dispute struck only the valley whose structure of trade
  ran to Tibet. The same event had no effect on the Ravi herders, whose trade
  ran south \cit{Kanum survey, PDF pp.~16, 51; Brahmaur survey, PDF p.~59}.
  The structure decided where the event was felt.
\item The sheep disease is an event inside a conjuncture. Braudel warns that
  \enquote{The short term is the most capricious, the most deceptive of time
  periods} \cit{\Braudel, PDF p.~6}. The fee-paying flocks did not recover
  after the disease \cit{Appendix Table~\ref{tab:series}}, and Newell
  attributes the loss to \enquote{the Forest Department permit system}
  \cit{Newell, PDF p.~37}.
\end{itemize}

\subsection{The Ravi and the Sutlej}
\begin{itemize}
\item The two valleys share the mountain and the season and differ in the
  direction of their rivers and in how the pastures were held \cit{Newell, PDF
  p.~18; Kanum survey, PDF p.~133}.
\item On the Ravi, where the raja had farmed out the grazing fee, the
  Department's permits became a market. On the Sutlej, where the pastures were
  commons, they became state forest under the same Department
  \cit{Newell, PDF pp.~18, 37; Kanum survey, PDF p.~133}.
\item The form of property before 1948 shaped the form after it. This is the
  dialectic at the scale of a valley: each new arrangement is built from the
  one before.
\item On the Sutlej the herders were also traders. Habib notes that mobile
  groups carried long-distance trade in Mughal India: \enquote{the famous nomadic
  class of banjaras} \cit{\Habib, PDF p.~41}. The Kanum flocks laden with
  goods for Lavi are a hill version \cit{Kanum survey, PDF p.~16}.
\end{itemize}

\subsection{Insights to Draw Out}
\begin{itemize}
\item The form of the migration lasted, and the people who gained from it
  changed. The months and routes stayed, and the claim on the product moved
  from the raja and his contractor to the Forest Department and the market.
\item Altitude was a condition for large flocks, and ownership decided who used
  it \cit{Appendix Tables~\ref{tab:alt} and~\ref{tab:class}}.
\item The state took the middleman's place, charged per permit holder instead
  of per district, and limited the number of animals. The limit created a
  market in permits \cit{Newell, PDF pp.~18, 19, 37; Brahmaur survey, PDF
  p.~46}.
\item Payment moved from kind to cash. The pohal was paid in meals, wool and
  animals, the village shepherd in grain, and the road labourer in cash
  \cit{Brahmaur survey, PDF pp.~45, 55; Tandi survey, PDF p.~65}. Habib on
  cash revenue: the payer \enquote{was compelled to sell his produce in order to
  get money to pay for it} \cit{\Habib, PDF p.~11}.
\item Where a household owned its means of production, its surplus was still
  taken through credit and the market \cit{Maingal survey, PDF pp.~42, 60}.
\item The sources speak through surveyors, settlement officers and headmen. The
  Lambardar of Nachar was the largest flock owner \cit{Nachar survey, PDF
  p.~33}, as Marx's \enquote{potail} was \enquote{collector or renter of the
  village} \cit{Marx, \enquote{The British Rule in India}, PDF p.~3}.
\item The Sipis, the Reharas, the Kolis, the hired fawals and the indebted
  Gujjars have no voice in the sources. Braudel: \enquote{There are others who
  matter but who are silent} \cit{\Braudel, PDF p.~18}.
\end{itemize}

% ===================================================================
\section{Conclusion}
\budget{200}

\begin{itemize}
\item Answer the question in two halves. The cycle persisted as an adaptation
  to mountain and season. Who could take part in it, and who gained from it,
  changed with each conjuncture.
\item By 1961 the migration ran on fixed dhars, under permits that could be
  bought and sold, with fewer fee-paying animals and with hired shepherds
  leaving for the roads \cit{Brahmaur survey, PDF p.~55; Newell, PDF pp.~37,
  38; Kanum survey, PDF p.~51}.
\item Those without permits, the castes that owned few animals and the indebted
  Gujjars lost by the change. Permit holders and the state gained.
\item Read alone, Braudel gives the frame and Marx gives the motion. Read
  together, the levels of time show where to look, and the material basis shows
  why things moved. State this once, in plain words.
\item The evidence graph draws almost all its claims from the 1961 volumes, so
  the colonial conjunctures are thinly attested \cit{Appendix
  Table~\ref{tab:graph}}.
\item Limits to name. A count of claims measures how often the sources mention a
  relation, not its size. The villages were chosen by the Census. The livestock
  and fee figures were read from page images. No source dates the closing of
  Shipki La, and no source maps the migration routes.
\end{itemize}

% ===================================================================
\appendix
\section{Data Tables}\label{app:data}

\noindent All values are computed by \path{scripts/build_evidence.py} and
appear in \path{essay/evidence/evidence.pdf}. Livestock and fee figures were
read from page images and have not been checked by a person.

\begin{table}[h]\centering\small
\caption{Pipeline counts.}\label{tab:pipe}
\begin{tabular}{lr}\toprule
Quantity & Count\\\midrule
Pages of the 1961 volumes scored / selected & 3,754 / 428\\
Pages of the 1881, Rose and 1931 volumes scored / selected & 2,939 / 82\\
Claims proposed by the language model & 1,843\\
Claims rejected, quotation not on the cited page & 44\\
Claims from image-read tables / entered by hand & 174 / 38\\
Claims kept in the evidence graph & 1,964\\
Second check: supported / partly / not supported / unclear & 1,135 / 545 / 108 / 1\\
Copied passages between sources & 4,698\\
Newell passages copied from Rose (all on PDF pp.~52 to 99) & 91\\
\bottomrule\end{tabular}
\end{table}

\begin{table}[h]\centering\small
\caption{Altitude and sheep and goats per household (selected villages). Spearman rank correlation of altitude with animals per household 0.54 across 19 villages, and with the landholding share −0.05 (project analysis A18).}\label{tab:alt}
\begin{tabular}{llrrrr}\toprule
Village & River & Altitude (m) & Households & Sheep and goats & Per household\\\midrule
Tandi & Chenab & 3048 & 18 & 384 & 21.3\\
Kanum & Sutlej & 2926 & 113 & 1827 & 16.2\\
Kothi & Sutlej & 2515 & 128 & 2557 & 20.0\\
Chaupal & Yamuna & 2342 & 80 & 126 & 1.6\\
Maingal & Ravi & 2286 & 31 & 405 & 13.1\\
Devi-Kothi & Ravi & 2164 & 60 & 1337 & 22.3\\
Brahmaur & Ravi & 2134 & 99 & 4484 & 45.3\\
Chitrari & Ravi & 1829 & 80 & 159 & 2.0\\
Bir & Beas & 1219 & 36 & 208 & 5.8\\
Batal & & 1088 & 140 & 40 & 0.3\\
Shakrori & Sutlej & 518 & 60 & 76 & 1.3\\
\bottomrule\end{tabular}
\end{table}

\begin{table}[h]\centering\small
\caption{Movements of the Gaddis and Gujjars named with months.}\label{tab:cal}
\begin{tabularx}{\textwidth}{llYY}\toprule
Group & Movement to & Years of the sources & Months named (years)\\\midrule
Gaddi & winter grounds & 1881, 1911, 1961 & Oct to Feb (1911, 1961)\\
Gaddi & summer pastures & 1911, 1961 & Mar to Nov (1961)\\
Gujjar & summer pastures & 1881, 1911, 1961 & Apr, May (1961)\\
\bottomrule\end{tabularx}
\end{table}

\begin{table}[h]\centering\small
\caption{Claims in the evidence graph by conjuncture and year of the source.}\label{tab:graph}
\begin{tabular}{lrrrrrr}\toprule
Conjuncture & 1881 & 1911 & 1914 & 1919 & 1931 & 1961\\\midrule
cycle & 4 & 11 & 2 & 6 & 2 & 344\\
two sovereigns & 3 & 1 & 0 & 0 & 0 & 1\\
raja's pastures & 8 & 12 & 9 & 5 & 0 & 72\\
Forest Department & 0 & 0 & 0 & 0 & 0 & 55\\
roads and wages & 11 & 7 & 2 & 0 & 0 & 169\\
trade and credit & 7 & 7 & 2 & 2 & 0 & 375\\
standing social order & 10 & 5 & 0 & 0 & 0 & 724\\
\bottomrule\end{tabular}

\smallskip
Of the 106 claims on the raja's pastures, 54 come from Newell (1961).
\end{table}

\begin{table}[h]\centering\small
\caption{Most frequent relations of payment in the graph, by conjuncture (claim counts).}\label{tab:flows}
\begin{tabular}{lrrr}\toprule
Who gives to whom (kind) & Raja's pastures & Forest Dept. & Trade and credit\\\midrule
households to moneylender (interest) & 0 & 0 & 42\\
households to co-operative society (interest) & 0 & 0 & 18\\
households to shopkeeper (produce sold) & 0 & 0 & 10\\
households to shopkeeper (interest) & 0 & 0 & 8\\
grazier to Forest Department (dues) & 0 & 8 & 0\\
shepherd to Forest Department (fee or tax) & 0 & 8 & 0\\
shepherd to Forest Department (permit) & 0 & 5 & 0\\
households to Raja of Chamba (dues) & 4 & 0 & 0\\
\bottomrule\end{tabular}

\smallskip
Main receivers by claims received: Raja of Chamba 11 and fee contractor 4
(raja's pastures), Forest Department 35, moneylender 53 and shopkeeper 29
(trade and credit).
\end{table}

\begin{table}[h]\centering\small
\caption{Keyword rates per 10,000 words in the colonial layers (whole volumes).}\label{tab:layers}
\begin{tabular}{lrrrrrrr}\toprule
Layer & Flocks & Pasture & Rights and dues & Forest and state & Trade & Debt & Wage work\\\midrule
1881 census & 1.82 & 1.43 & 0.31 & 0.00 & 3.34 & 0.97 & 0.87\\
Rose (1911) & 7.61 & 1.85 & 1.35 & 0.00 & 4.71 & 1.23 & 0.52\\
1931 census & 0.51 & 0.29 & 1.28 & 0.04 & 1.19 & 1.93 & 2.42\\
\bottomrule\end{tabular}
\end{table}

\begin{table}[h]\centering\small
\caption{Sheep and goats by caste class, all caste tables.}\label{tab:class}
\begin{tabular}{lrr}\toprule
Caste class & Sheep and goats & Share (\%)\\\midrule
landholding & 10,902 & 89.3\\
scheduled caste & 992 & 8.1\\
herding tribe & 190 & 1.6\\
other & 106 & 0.9\\
trading & 20 & 0.2\\
all caste tables & 12,210 & 100.0\\
\bottomrule\end{tabular}
\end{table}

\begin{table}[h]\centering\small
\caption{Landholding share by village (selected).}\label{tab:own}
\begin{tabular}{llrrl}\toprule
Village & River & Sheep and goats & Landholding (\%) & Largest owning group\\\midrule
Brahmaur & Ravi & 4484 & 97.9 & Rajput\\
Kothi & Sutlej & 2557 & 94.5 & Rajput\\
Kanum & Sutlej & 1827 & 95.5 & Kanet (Rajput)\\
Kupha & Chenab & 596 & 93.0 & Rajput\\
Tandi & Chenab & 384 & 92.2 & Bodh\\
Maingal & Ravi & 405 & 54.6 & Brahmin\\
Rajana & Yamuna & 559 & 40.1 & Koli\\
Pangna & & 62 & 4.8 & Koli\\
Chaunri & Yamuna & 45 & 8.9 & Koli\\
\bottomrule\end{tabular}
\end{table}

\begin{table}[h]\centering\small
\caption{The Brahmaur livestock table (Brahmaur survey, PDF p.~54).}\label{tab:brahmaur}
\begin{tabular}{lrrr}\toprule
Caste & Households & Owning sheep or goats & Sheep and goats\\\midrule
Brahman & 34 & 20 & 2027\\
Rajput & 49 & 20 & 2364\\
Sipi & 15 & 8 & 93\\
Rehara & 1 & 0 & 0\\
total & 99 & 48 & 4484\\
\bottomrule\end{tabular}
\end{table}

\begin{table}[h]\centering\small
\caption{Trade, flock and debt words per 10,000 words of each survey.}\label{tab:trade}
\begin{tabular}{llrrr}\toprule
Survey & River & Trade & Flocks & Debt\\\midrule
Kanum & Sutlej & 33.6 & 17.81 & 8.03\\
Kothi & Sutlej & 14.33 & 17.82 & 10.1\\
Nachar & Sutlej & 13.0 & 16.02 & 10.45\\
Devi-Kothi & Ravi & 9.17 & 14.74 & 6.88\\
Chitrari & Ravi & 8.43 & 20.16 & 5.96\\
Brahmaur & Ravi & 8.29 & 27.09 & 4.67\\
Maingal & Ravi & 6.62 & 20.53 & 17.55\\
\bottomrule\end{tabular}
\end{table}

\begin{table}[h]\centering\small
\caption{Grazing charges. The units differ by row, so compare only rows with the same unit.}\label{tab:fees}
\begin{tabularx}{\textwidth}{llllYr}\toprule
Year & Place & Collector & Animal & Rate as printed, and its unit & Rupees\\\midrule
1950 & Chamba State & fee contractor & goat & 2 annas 6 pies per head per grazing district & 0.156\\
1950 & Chamba State & fee contractor & sheep & 1 anna 6 pies per head per grazing district & 0.094\\
1961 & Brahmaur & Forest Department & sheep or goat & Rs 15 per hundred heads per summer season, paid twice a year & 0.15 per head\\
1962 & Devi-Kothi & Forest Department & sheep & Rs 0.10 per head, permit valid 1 June to 15 October (both classes) & 0.10\\
1962 & Devi-Kothi & Forest Department & goat & Rs 0.25 per head (both classes) & 0.25\\
1962 & Devi-Kothi & Forest Department & cow or bull & Rs 1.00 per head (both classes) & 1.00\\
1962 & Devi-Kothi & Forest Department & buffalo & Rs 6.00 non-right holders, Rs 4.50 right holders & 6.00 / 4.50\\
1962 & Lahul & Forest Department & goat & 4 annas per head, period not stated & 0.25\\
1962 & Lahul & Forest Department & sheep & 2 annas per head, period not stated & 0.125\\
1962 & Brahmaur & Forest Department & sheep or goat & 2 annas per head, period not stated & 0.125\\
1962 & Kangra & Forest Department & sheep & 8 naye paise per head, period not stated & 0.08\\
1962 & Kangra & Forest Department & goat & 4 annas per head, period not stated & 0.25\\
\bottomrule\end{tabularx}

\smallskip
Sources are Newell, PDF p.~18, the Brahmaur survey, PDF p.~55, and the
Devi-Kothi survey, PDF p.~47 (Table XIV, read from the page image). Penalties
for grazing without a permit in Devi-Kothi are Rs 9.00 and Rs 6.00 for a
buffalo, Rs 0.50 for a sheep and Rs 1.00 for a goat. The 1962 rates for
Lahul, Brahmaur and Kangra are from Newell, PDF p.~19, read from the page
image because the OCR is unreadable. Newell's Brahmaur rate, 2 annas
(Rs 0.125) per head differs from the survey's Rs 15 per hundred heads
(Rs 0.15) per summer season, and Newell does not state the period.
\end{table}

\begin{table}[h]\centering\small
\caption{Fee-paying sheep and goats in Brahmaur tahsil (Newell, PDF p.~38), and the change by phase.}\label{tab:series}
\begin{tabular}{lrrr}\toprule
Year & Sheep & Goats & Total\\\midrule
1952--53 & 42,176 & 31,695 & 73,871\\
1956--57 & 40,600 & 30,896 & 71,496\\
1959--60 & 30,828 & 24,864 & 55,692\\
1961--62 & 30,515 & 24,406 & 54,921\\\midrule
\multicolumn{3}{l}{Change 1952--53 to 1956--57 (\%)} & −3.2\\
\multicolumn{3}{l}{Change 1956--57 to 1959--60 (\%)} & −22.1\\
\multicolumn{3}{l}{Change 1959--60 to 1961--62 (\%)} & −1.4\\
\bottomrule\end{tabular}
\end{table}

\begin{table}[h]\centering\small
\caption{Units of money, grain and height in the sources.}\label{tab:units}
\begin{tabularx}{\textwidth}{lY}\toprule
Unit & Value\\\midrule
rupee, anna, pie & 1 rupee = 16 annas, and 1 anna = 12 pies. So 2 annas 6 pies = Rs 0.156.\\
naya paisa & From 1957, 1 rupee = 100 naye paise. Sources of 1961 and 1962 mix both systems.\\
maund, seer & 1 maund = 40 seers. A grain loan at 5 seers per maund a year is 12.5\% interest.\\
patha & About one seer \cit{Tandi survey, PDF p.~65}.\\
foot, metre & The surveys state heights in feet. The appendix tables give metres (1 foot = 0.3048 m).\\
\bottomrule\end{tabularx}
\end{table}

\begin{table}[h]\centering\small
\caption{Wages, prices and interest, with their units.}\label{tab:wages}
\begin{tabularx}{\textwidth}{lYl}\toprule
Item & Amount and unit & Source\\\midrule
Pohal (hired shepherd), Brahmaur & Rs 2 per month in cash, with free meals, clothing, wool, and two or four sheep per hundred animals & Brahmaur, p.~55\\
Fawal (hired shepherd), Kanum & Rs 60 per year each, with meals, 8 kg of wool, shoes and clothes & Kanum, p.~132\\
Village shepherd, Tandi & one patha of buckwheat per sheep or goat per season, from each household & Tandi, p.~65\\
Grazier, Devi-Kothi & 20 to 25 naye paise per sheep or goat, May or June to October & Devi-Kothi, p.~48\\
Road and timber labour, Brahmaur & about Rs 2 per day & Brahmaur, p.~45\\
Timber labour, Devi-Kothi & about Rs 4 per day & Devi-Kothi, p.~48\\
Gujjar loans, Maingal & 30\% interest per year & Maingal, p.~42\\
Grain loans, Kanum & 5 seers per maund per year (12.5\%) & Kanum, p.~55\\
Ghee, Devi-Kothi markets & Rs 12 per batti (Chamba tole), 1962 & Devi-Kothi, p.~47\\
\bottomrule\end{tabularx}

\smallskip
A pohal paid by the month and a road labourer paid by the day cannot be
compared on the cash alone, because the pohal's pay was mostly in kind.
\end{table}

\clearpage
\section{References}

\subsection*{Primary Sources}
\begin{itemize}
\item Ibbetson, Denzil. \emph{Report on the Census of the Panjab, 1881}, Results volume. Calcutta, 1883.
\item Rose, H.\,A. \emph{A Glossary of the Tribes and Castes of the Punjab and North-West Frontier Province}, vols.~1 to~3. Lahore, 1911 to 1919.
\item Khan, Khan Ahmad Hasan. \emph{Census of India, 1931, Vol.~XVII: Punjab}, Part~I, Report. Lahore, 1933.
\item Census of India 1961, Vol.~XX (Himachal Pradesh), Part~VI, Village Survey Monographs: Batal, Bir, Brahmaur, Chaunri, Chaupal, Chergaon, Chitrari, Devi-Kothi, Kanum, Kothi, Kupha (Pangi), Maingal, Nachar, Pangna, Rajana, Shakrori.
\item Census of India 1961, Vol.~XIII (Punjab), Part~VI, Village Survey Monograph: Tandi.
\item Census of India 1961, Vol.~XIII, Part~IX, \emph{Punjab Census Atlas}.
\item Newell, William H. \enquote{The Gaddi and Affiliated Castes of the Western Himalayas.} In Census of India 1961, \emph{Report on Scheduled Castes and Scheduled Tribes of Himachal Pradesh}.
\item Census of India 1961, Himachal Pradesh, \emph{District Census Handbook: Mahasu}, and \emph{Craft Survey: Weaving}.
\end{itemize}

\subsection*{Course Readings}
\begin{itemize}
\item Alam, Aniket. \enquote{Marxist History.} Lectures 8 to 10, Introduction to History (HS3.201), IIIT Hyderabad, Monsoon 2026.
\item Aymard, Maurice. \enquote{The Cost of War.} In \emph{French Studies in History}, vol.~2, \emph{Departures}, 418--427.
\item Bloch, Marc. \emph{The Historian's Craft}. Translated by Peter Putnam. New York: Vintage, 1953. Chapter~III.
\item Braudel, Fernand. \enquote{History and the Social Sciences: The Longue Durée.} Translated by Immanuel Wallerstein. \emph{Review (Fernand Braudel Center)} 32, no.~2 (2009): 171--203.
\item Habib, Irfan. \enquote{Potentialities of Capitalistic Development in the Economy of Mughal India.} \emph{Journal of Economic History} 29, no.~1 (1969): 32--78.
\item Kosambi, D.\,D. \enquote{Social and Economic Aspects of the Bhagavad Gita.} \emph{Journal of the Economic and Social History of the Orient} 4, no.~2 (1961): 198--224.
\item Le Roy Ladurie, Emmanuel. \enquote{The History of Rain and Fine Weather.} In \emph{French Studies in History}, vol.~1, \emph{The Inheritance}, 192--214.
\item Marx, Karl. \enquote{The British Rule in India.} \emph{New-York Daily Tribune}, 25 June 1853.
\item Marx, Karl. \enquote{The Future Results of British Rule in India.} \emph{New-York Daily Tribune}, 8 August 1853.
\item Marx, Karl, and Friedrich Engels. \emph{Manifesto of the Communist Party}. Excerpt, pp.~14--21.
\end{itemize}

\end{document}
